%% ch07 --- Importy i wstrzykiwanie zależności: moduły, cykle i korzeń
%% kompozycji.
%%
%% Napisany. Bliźniak angielski musi mieć tę samą strukturę, te same liczby w
%% tej samej kolejności i te same klucze listingów, rysunków i ćwiczeń;
%% sprawdzenia C4, C12 i C14 porównują wszystkie trzy.
%%
%% Decyzja tłumaczeniowa: lifetime to "czas życia", nie "cykl życia" --
%% "cykl" w tym rozdziale znaczy cykl importów i kolizja byłaby niepotrzebna.

\chapter{Importy i wstrzykiwanie zależności: moduły, cykle i korzeń kompozycji}
\label{ch:imports}

\noindent
Dyrektywa \code{using} w C\# jest instrukcją dla kompilatora, gdzie szukać
nazw. Niczego nie uruchamia. Nic w wskazanym przez nią zestawie nie dzieje
się dlatego, że ją napisałeś, kolejność dyrektyw na górze pliku nie ma
znaczenia, a napisanie jednej dwa razy jest kwestią stylu, nie zachowania.

Nic z tego nie jest prawdą w Pythonie i prawie wszystko, co psuje się przy
importach, wynika z jednego zdania, które to zastępuje:
\textbf{moduł jest obiektem, zbudowanym przez wykonanie pliku, raz}.
\code{import} jest wyrażeniem, które go buduje. Ten rozdział traktuje to
zdanie poważnie przez pięć sekcji, a na końcu okazuje się, że rozstrzygnęło
pytanie, które wygląda na niezwiązane: gdzie podziewa się wstrzykiwanie
zależności, gdy język nie ma kontenera.

\begin{csbox}
Najbliższą konstrukcją C\# nie jest wcale \code{using}. Jest nią klasa
statyczna ze statycznym konstruktorem: jedna instancja na proces,
zainicjowana przez wykonanie kodu, przy pierwszym kontakcie. Wersja
Pythona różni się dwiema rzeczami, które mają znaczenie. Wyzwala ją
\emph{import}, a nie pierwsze użycie, i język nie daje o niej obietnicy
bezpieczeństwa wątkowego, na której warto polegać.
\end{csbox}

\section{Moduł jest obiektem, który wykonuje się raz}
\label{sec:module-object}
\index{moduły!jako obiekty}

Oto moduł, którego ciało robi coś, co widać.

\pyfile{code/ch07/greeting.py}{Ciało modułu to kod i on się wykonuje}{lst:ch07-greeting}

\noindent
Teraz drugi plik importuje go trzy razy. Zanim to uruchomisz: ile razy ciało
coś wypisze i czy te trzy nazwy to ten sam obiekt?

\pyfile{code/ch07/runs_once.py}{Trzy importy jednego modułu}{lst:ch07-runs-once}

\transcript{ch07-runs-once}

\noindent
Raz, i tak. Mechanizmem jest słownik. \code{import x} szuka klucza
\code{"x"} w \api{sys.modules}; pudło wykonuje plik od góry do dołu i wstawia
powstały obiekt modułu do cache pod tym kluczem, a trafienie zwraca to, co
już tam jest, i nie wykonuje niczego. \api{importlib.import\_module} to ta
sama operacja zapisana jako funkcja, i dlatego listing może nią zastąpić
drugiego i trzeciego importującego.

\mermaidfig{ch07-import-once}{Co robi \code{import x}. Ciało \code{x} wykonuje się przy pudle i nigdy więcej}{instrukcja importu jako zajrzenie do cache}

\noindent
Wynikają z tego natychmiast trzy konsekwencje i wszystkie trzy inżynierowie
przychodzący z .NET rozumieją opacznie.

\begin{itemize}[leftmargin=1.4em, itemsep=3pt]
  \item \textbf{Import nie jest darmowy.} Cokolwiek robi ciało \dash{}
        czyta plik, otwiera połączenie, zadaje pytanie środowisku \dash{}
        robi to przy czyimś pierwszym imporcie, którym zwykle jest
        najwolniejsza komenda, jaką masz, a nie ta, która tego
        potrzebowała.
  \item \textbf{Nazwa na poziomie modułu jest globalna dla procesu.} Obiekt
        jest jeden, więc rzecz mutowalna związana na poziomie modułu jest
        dzielona przez wszystkich importujących. To pole statyczne, bez
        żadnej ceremonii, która kazałaby je uznać za decyzję.
  \item \textbf{Kolejność importów staje się obserwowalna.}
        \code{SALUTATION} z listingu~\ref{lst:ch07-greeting} czyta
        środowisko w czasie importu, więc to, co trzyma, zależy od tego,
        kiedy nastąpił pierwszy import, a test, który ustawi zmienną
        później, nie zmieni niczego.
\end{itemize}

\begin{trapbox}
\textbf{\enquote{Moduł to przestrzeń nazw; zaimportowanie go jest darmowe i
czyste.}} Jest obiektem, a zbudowanie go wykonuje plik. Efekt uboczny na
poziomie modułu zdarza się raz na cały proces, w momencie wyznaczonym przez
kolejność importów, a nie przez ciebie \dash{} i najtańszą poprawką jest
prawie zawsze przeniesienie pracy do funkcji, gdzie o momencie decyduje
wołający.
\end{trapbox}

\begin{exercise}{e07_01_import_time_work}{Wynieś pracę poza czas importu}
Plik startowy buduje słownik ustawień ze środowiska na poziomie modułu.
Usuń go i daj modułowi funkcję \code{build\_settings()}, która wykona tę
samą pracę, gdy zostanie wywołana. Jeden z testów importuje moduł, zmienia
środowisko i \emph{dopiero potem} woła funkcję: wyrażenie słownikowe na
poziomie modułu nie ma szans go przejść i o to właśnie chodzi.
\end{exercise}

\section{Pakiety i dwa sposoby uruchomienia pliku}
\label{sec:packages}
\index{moduły!pakiety}

Pakiet to katalog z plikiem \code{\_\_init\_\_.py}, a użytecznie jest myśleć
o nim tak, że pakiet też jest modułem: \code{\_\_init\_\_.py} \emph{jest}
ciałem pakietu, wykonuje się raz, zanim cokolwiek w środku, a to, co zwiąże,
jest tym, co \code{from shop import x} może znaleźć. Plik
\code{\_\_init\_\_.py}, który dla wygody wołających importuje własne
podmoduły, każe każdemu importującemu którykolwiek z nich zapłacić za
wszystkie \dash{} i dlatego domyślnie w tej książce ten plik jest pusty.

Zostają dwa pytania, które czytelnik naprawdę już spotkał: importy
bezwzględne kontra względne oraz po co istnieje \code{python -m}. Mają jedną
odpowiedź. Ten moduł melduje, za co interpreter go uważa.

\pyregion{code/ch07/shop/where.py}{relative}{Import względny, próbowany w czasie importu}{lst:ch07-relative}

\noindent
A ten uruchamia ten sam plik na oba sposoby.

\pyfile{code/ch07/two_ways.py}{Ten sam plik, uruchomiony jako plik i jako moduł}{lst:ch07-two-ways}

\transcript{ch07-two-ways}

\noindent
Uruchomiony jako plik, moduł nie ma pakietu \dash{} \code{\_\_package\_\_}
to \code{None} \dash{} więc import względny nie ma \emph{względem} czego
być, a porażką jest \code{ImportError} mówiący o pakiecie nadrzędnym, a nie
o brakującym pliku. Uruchomiony przez \code{-m}, interpreter importuje go
jako część pakietu \code{shop} i import względny się rozwiązuje. Zwróć
uwagę na drugą różnicę, bo to ona gryzie później: na czele \api{sys.path}
stoi \emph{katalog samego skryptu} w pierwszym przypadku, a katalog, z
którego uruchamiasz, w drugim.

\begin{csbox}
\code{-m} jest najbliższe uruchomieniu projektu, a nie pliku zestawu:
\code{dotnet run} rozwiązuje projekt i jego referencje, podczas gdy
wywołanie biblioteki DLL wprost daje to, co akurat leży w katalogu. Układ
\code{src/} z rozdziału~\ref{ch:packaging} istnieje po to, by ten drugi
rodzaj wypadku był niemożliwy, a ta sekcja jest mechanizmem pod nim.
\end{csbox}

\noindent
Ten wpis w \api{sys.path} jest \emph{pierwszy}, więc twój plik wygrywa z
czymkolwiek o tej samej nazwie niżej, z biblioteką standardową włącznie.
Cztery linie poniżej nie są błędne. Skopiuj je do katalogu pod jedną nazwą i
działają; pod inną nie.

\pyfile{code/ch07/shadow/innocent.py}{Cztery linie, które nie zależą od nazwy swojego pliku}{lst:ch07-innocent}

\pyfile{code/ch07/shadowing.py}{Te same cztery linie, pod dwiema nazwami}{lst:ch07-shadowing}

\transcript{ch07-shadowing}

\noindent
\code{import random} w pliku nazwanym \code{random.py} znajduje ten właśnie
plik, wiąże moduł z nazwą, a moduł nie ma na sobie funkcji \code{random}.
Komunikat nie nazywa ani przesłonięcia, ani pliku, który je spowodował.

\begin{trapbox}
\textbf{\enquote{Mogę nazwać plik albo zmienną \code{list}, \code{id},
\code{type} czy \code{random}.}} Na poziomie modułu nazwa wygrywa z
biblioteką standardową dla wszystkiego niżej w \api{sys.path}; wewnątrz
funkcji wygrywa z nazwą wbudowaną do końca zasięgu. Oba psują się w pewnej
odległości od linii, która je spowodowała, a komunikat nazywa objaw.
\end{trapbox}

\noindent
Ostatnim kawałkiem publicznej powierzchni modułu jest ten, który C\# zapisuje
modyfikatorem dostępu. Python nie ma \code{internal}: wiodące podkreślenie
jest konwencją, a \code{\_\_all\_\_} jest jedyną rzeczą, która nadaje
\code{from module import *} określone znaczenie. Bez niego gwiazdka kopiuje
każdą nazwę niezaczynającą się od podkreślenia, w tym moduły, które ten plik
zaimportował \dash{} więc import z gwiazdką potrafi wręczyć ci cudzy
\code{os}.

\begin{exercise}{e07_02_public_surface}{Powiedz, co moduł eksportuje}
Daj plikowi startowemu \code{\_\_all\_\_} wymieniające tylko te dwie
funkcje, których powinni używać wołający z zewnątrz. Testy sprawdzają, co
skopiowałby import z gwiazdką, a także to, że pomocnicze funkcje wciąż są
dostępne po nazwie, bo \code{\_\_all\_\_} rządzi gwiazdką i niczego nie
ukrywa.
\end{exercise}

\begin{trapbox}
\textbf{\enquote{\code{from x import *} to \code{using x;}.}} Kopiuje nazwy
do twojego modułu, zamiast czynić je widocznymi, więc na pytanie, skąd
wzięła się dana nazwa, nie da się już odpowiedzieć czytaniem pliku. Wybieraj
\code{import x} albo \code{from x import name} dla garstki nazw, których
używasz bez przerwy.
\end{trapbox}

\section{Cykliczny import to nie ten błąd, o którym myślisz}
\label{sec:cycles}
\index{moduły!importy cykliczne}
\index{diagnostyka!import cykliczny}

Dwa moduły importujące się nawzajem to kształt, którego wszystkim kazano się
bać. Zanim przeczytasz dalej, rozstrzygnij, która z tych par zawiedzie: ta,
w której każdy robi \code{import} drugiego, czy ta, w której każdy robi
\code{from} drugiego \code{import} nazwy.

Pierwsza działa. Sam cykl nigdy nie jest błędem, bo do
\api{sys.modules} trafia \emph{obiekt} modułu, zanim wykona się jego ciało
\dash{} taka jest kolejność i to ona zatrzymuje rekurencję. Zawodzi proszenie
półwykonanego modułu o coś, czego jeszcze nie związał, a \code{from x import
y} prosi dokładnie o to, w czasie importu.

\pyfile{code/ch07/cycles.py}{Ten sam cykl, na cztery sposoby}{lst:ch07-cycles}

\transcript{ch07-cycles}

\mermaidfig{ch07-cycle-break}{Gdzie cykl pęka, a gdzie nie: obiekt modułu wiąże się od razu, nazwa w nim nie}{cykl i linia, na której zawodzi}

\noindent
Wszystkie trzy wyjścia są więc sposobami na przesunięcie wyszukania poza
czas importu i tylko ostatnie z nich usuwa cykl.

\begin{enumerate}[leftmargin=1.6em, itemsep=3pt]
  \item \textbf{Trzymaj moduł, a nie nazwę.} \code{import x}, a potem
        \code{x.NAME} tam, gdzie jest używane. Związanie modułu udaje się
        zawsze; zanim cokolwiek odczyta z niego atrybut, ciało już się
        skończy.
  \item \textbf{Importuj wewnątrz funkcji.} Po pierwszym wywołaniu to
        zajrzenie do słownika, a całe pytanie wychodzi poza kolejność
        importów. Kosztem jest to, że czytelnik musi tego poszukać, więc
        umieść to w najmniejszym zasięgu, który działa.
  \item \textbf{Wynieś do trzeciego modułu to, czego chciały obie strony.}
        To wyjście jest zwykle właściwe, a cykl zwykle jest dowodem, że było
        należne: dwa moduły, które siebie potrzebują, na ogół dzielą coś
        trzeciego, czego żaden z nich nie jest właścicielem.
\end{enumerate}

\begin{warning}
\textbf{Diagnoza zależy od tego, gdzie leży plik, a nie od tego, co poszło
źle.} Dwa ostatnie wpisy w powyższym zapisie to ta sama usterka. Wewnątrz
pakietu CPython mówi \emph{partially initialized module \dots{} most likely
due to a circular import}, co nazywa przyczynę. W dwóch plikach obok
uruchamianego skryptu mówi \emph{consider renaming \dots{} if it has the
same name as a library you intended to import}, co dotyczy pułapki z
\S\ref{sec:packages}, a nie tej. Oba zostały uruchomione, żeby ten zapis
powstał. Drugi komunikat pojawia się dokładnie w tym układzie, który
przyjmuje pierwszy projekt w Pythonie, więc jeśli kiedyś goniłeś kolizję
nazw, której nie było, to właśnie dlatego.
\end{warning}

\begin{trapbox}
\textbf{\enquote{Import cykliczny to błąd.}} To fakt o twoim grafie modułów
i Python go toleruje. Błędem jest \emph{nazwa} czytana z modułu, którego
ciało do niej nie doszło, i dlatego zamiana \code{from x import y} na
\code{import x} czasem naprawia build bez żadnej innej zmiany \dash{} i
dlatego ta poprawka, wzięta sama, zostawia cykl na miejscu.
\end{trapbox}

\begin{exercise}{e07_03_defer_the_import}{Płać za import wtedy, gdy jest używany}
Plik startowy importuje moduł z biblioteki standardowej na górze i używa go
w jednej funkcji, do której większość wołających nigdy nie dociera. Przenieś
import do wnętrza tej funkcji. Testy czyszczą cache, importują moduł i
sprawdzają, że nic nie przyszło, a potem wołają funkcję i sprawdzają, że
przyszło.
\end{exercise}

\section{Korzeń kompozycji jest funkcją}
\label{sec:composition-root}
\index{wstrzykiwanie zależności!korzeń kompozycji}

Nic z poprzednich trzech sekcji nie dotyczy wstrzykiwania zależności i
wszystko z nich dotyczy. Kontener istnieje po to, by obiekty nie sięgały
same po swoich współpracowników; w Pythonie sięgają po nazwę na poziomie
modułu, o której \S\ref{sec:module-object} pokazała, że jest globalna dla
procesu i wiązana w niewygodnym momencie. Zabierz to, a temu, co zostanie,
nie będzie potrzebna żadna maszyneria.

Zacznij od tego, czego potrzebuje praca, wyrażonego jako \code{Protocol}:
strukturalnie, więc klasa spełnia go kształtem i nigdy tego nie deklaruje.

\pyregion{code/ch07/composition.py}{ports}{Dwa protokoły: czego potrzebuje usługa i nic ponadto}{lst:ch07-ports}

\noindent
Praca bierze swoich współpracowników jako argumenty, jak wszystko inne.

\pyregion{code/ch07/composition.py}{service}{Usługa nie nazywa żadnego konkretnego typu}{lst:ch07-service}

\noindent
Implementacje są zwykłymi klasami. Żadna nie dziedziczy po protokole i żadna
o nim nie wspomina.

\pyregion{code/ch07/composition.py}{adapters}{Prawdziwe implementacje, które niczego nie deklarują}{lst:ch07-adapters}

\noindent
A korzeń kompozycji to jedna funkcja: jedyne miejsce w programie, w którym
pada nazwa konkretnej klasy.

\pyregion{code/ch07/composition.py}{root}{Korzeń kompozycji, w całości}{lst:ch07-root}

\mermaidfig{ch07-composition-root}{Montaż bez kontenera: jedna funkcja na brzegu zna konkretne typy i nic poza nią}{korzeń kompozycji}

\noindent
\api{functools.partial} wiąże początkowe argumenty i zostawia resztę
sygnatury w spokoju, więc wraca usługa z już wykonanym montażem. To właśnie
wręcza kontener, złożone ręcznie, w jednej linii, którą da się przeczytać
\dash{} a ponieważ \code{build()} ma adnotację mówiącą, że zwraca zwykły
obiekt wywoływalny, nic dalej nie sięgnie przez nią i nie zapyta, czym to
naprawdę jest. Checker tego pilnuje.

\begin{csbox}
\code{IServiceCollection} jest rejestrem, a rejestracja jest drugim opisem
programu, który trzeba utrzymywać w prawdzie. Korzeń kompozycji jest
programem: jeśli brakuje współpracownika, wywołanie nie przechodzi przez
checker i nie przechodzi przez test. Nic nie jest rozwiązywane w czasie
działania, więc nie ma tu czegoś takiego jak kontener, który wstaje, a
potem rzuca wyjątkiem przy pierwszym żądaniu.
\end{csbox}

\begin{exercise}{e07_04_composition_root}{Napisz korzeń kompozycji}
Plik startowy ma protokoły, pracę i funkcję \code{build}, która rzuca
wyjątkiem. Napisz ją. Test montuje wtedy usługę z atrapą zegara i listą,
z których żadne po niczym nie dziedziczy, i sprawdza, że dwa osobne montaże
nie dzielą niczego.
\end{exercise}

\section{Czasy życia i kontener, który naprawdę spotkasz}
\label{sec:lifetimes}
\index{wstrzykiwanie zależności!czasy życia}

Trzy czasy życia to pierwsze, czego inżynierowi .NET zabraknie. Każdy ma
odpowiedź po stronie Pythona i żadna z nich nie jest rejestracją.

\begin{center}
\small
\begin{tabularx}{\linewidth}{@{}lX@{}}
\toprule
\textbf{C\#} & \textbf{Python} \\
\midrule
\code{AddSingleton} & Obiekt zbudowany raz w korzeniu kompozycji i
  przekazywany w dół. Obiekt na poziomie modułu to drugi sposób i ten, z
  którym trzeba uważać: powstaje w czasie importu, wedle reguł z
  \S\ref{sec:module-object}. \\
\code{AddScoped} & Parametr. Zasięgiem jest ta funkcja, która trzyma obiekt
  \dash{} obsługa żądania, zadanie, test \dash{} a blok \code{with} go
  zamyka. \\
\code{AddTransient} & Zawołaj fabrykę. \code{partial} wiąże fabrykę
  równie chętnie jak instancję. \\
\code{IServiceProvider} & Zwykle nic. Tam, gdzie framework go potrzebuje,
  przynosi własny. \\
\bottomrule
\end{tabularx}
\end{center}

\noindent
Zostaje uczciwe pytanie, o które prosił opis tego rozdziału: czy jest
kontener wart użycia? Wyszukane na PyPI, a nie odtworzone z pamięci, w
\pinnedon{}: odpowiedź brzmi, że kilka jest utrzymywanych i aktywnie
wydawanych \dash{} między innymi \pkg{dependency-injector}, \pkg{injector},
\pkg{punq}, \pkg{svcs}, \pkg{wireup} i \pkg{kink} \dash{} i że żaden z nich
nie jest tym, którego używają wszyscy. To jest ustalenie. W .NET kontener
przychodzi z platformą, więc jego użycie jest domyślne, a nie decyzją; w
Pythonie przyjęcie któregoś jest decyzją, której będziesz bronić przed
następną osobą, przeciwko korzeniowi kompozycji, który jest funkcją.

\begin{trapbox}
\textbf{\enquote{Potrzebuję kontenera DI.}} Potrzebujesz tego, po co
kupiono kontener, czyli tego, żeby obiekty nie sięgały same po swoich
współpracowników. Robi to funkcja i \api{functools.partial}, typuje to
\code{Protocol}, a test podstawia cokolwiek bez rejestrowania. Sięgnij po
kontener, kiedy umiesz nazwać, co ci tutaj kupi \dash{} i zauważ, że ten,
którego większość czytelników naprawdę spotka, wcale nie jest kontenerem
ogólnego przeznaczenia.
\end{trapbox}

\noindent
Jest nim \api{Depends} z FastAPI i jest to kontener o bardzo wąskim
zakresie: montuje argumenty funkcji obsługującej żądanie, na żądanie. Warto
go tu zapowiedzieć tylko po to, by powiedzieć, czym jest, bo wszystko, co
czyni go interesującym \dash{} jakie naprawdę ma czasy życia, co robi
zależność z \code{yield}, gdzie idzie stan o czasie życia aplikacji \dash{}
należy do rozdziału~\ref{ch:services}, w którym mieszka ten framework.

\begin{versionbox}
\api{Depends} w \pkg{fastapi}~\fastapiver{} bierze obiekt wywoływalny oraz
\code{use\_cache} i \code{scope}. Czytaj tę sygnaturę z zainstalowanego
pakietu, a nie z tutoriala, bo tutorial jest twierdzeniem o jakiejś wersji;
rozdział~\ref{ch:services} tak robi.
\end{versionbox}

\noindent
Zysk, który rozdział obiecał, da się teraz wyrazić jednym zdaniem i jest to
własność strukturalna, a nie kwestia stylu: \textbf{pakiet, w którym żaden
moduł nie zmienia niczego globalnego w czasie importu i w którym jedyną
funkcją nazywającą konkretną klasę jest ta na brzegu, to pakiet, w którym
test podstawi cokolwiek bez łatania}. Obie połowy należą do tego rozdziału.
Pierwsza jest tym, po co były \S\ref{sec:module-object} do
\S\ref{sec:cycles}; druga to \S\ref{sec:composition-root}, która ma
trzydzieści linii dlatego, że pierwsza połowa odwaliła robotę.
